The Lifted \strips (LS) Transformer follows the same conceptual design
as the \strips Transformer of~\citet{nunez2026strips}, but extends it
from propositional to lifted action models.

Consider a \strips model $M=\langle\mathcal{P},A\rangle$. The applicability
of an action $a_i$ depends on the truth of its preconditions, each determined
by the most recent preceding action that affects the corresponding ground
atom. Masked hard attention implements this recency-based dependency by
retrieving the rightmost such action. Queries identify preconditions of the
current action, keys identify effects of preceding actions on the same atom,
and values indicate whether these effects make the atom true or false.
This is the core idea behind the \strips Transformer. In the lifted setting,
identifying the same ground atom additionally requires matching both its
predicate $p$ and object arguments $\mathbf{o}$.

The LS Transformer consists of a single self-attention block with one head
for each predicate $p$, which tracks the truth values of its ground
atoms along the trace. To distinguish these atoms, the architecture encodes the 
possible occurrences of each predicate within an action schema through
\emph{action patterns}. For a predicate $p$ of arity $r$, an action pattern 
$\alpha[\rho]$, with $\rho=(\rho_1,\ldots,\rho_r)$, specifies how the arguments 
of $p$ bind to the variables of $\alpha$. For example, $\alpha[3,1]$ represents 
the occurrence $p(x_3,x_1)$ in $\alpha(x_1,x_2,x_3)$.

For each predicate $p\in\mathcal{P}$, action schema $\alpha\in A$, and valid
action pattern $\alpha[\rho]$, the LS Transformer has three scalar values
$q_{p,\alpha,\rho}$, $k_{p,\alpha,\rho}$, and $v_{p,\alpha,\rho}$ in
$[0,1]$, used in its query, key, and value computations. Their binary
values for a \strips model $M$ are
\begin{align}
q_M(p,\alpha,\rho)
&=
\mathbb{I}\!\left[
p(x_{\rho_1},\ldots,x_{\rho_r})\in\mathrm{pre}(\alpha)
\right],\\
k_M(p,\alpha,\rho)
&=
\mathbb{I}\!\left[
p(x_{\rho_1},\ldots,x_{\rho_r})
\in\mathrm{add}(\alpha)\cup\mathrm{del}(\alpha)
\right],\\
v_M(p,\alpha,\rho)
&=
\mathbb{I}\!\left[
p(x_{\rho_1},\ldots,x_{\rho_r})\in\mathrm{del}(\alpha)
\right].
\end{align}
Thus, $q_M$ indicates whether the predicate occurrence is a precondition,
$k_M$ whether it is an effect, and $v_M$ whether that effect is a delete
effect. During learning, $q$, $k$, and $v$ are differentiable relaxations
of these binary values; their parameterization is detailed in
Appendix~\ref{app:strips-transformer}.

The head for predicate $p$ compares an occurrence of $p$ in the current action
$a_i=\alpha_i(\mathbf{o}_i)$ with an occurrence in a preceding action
$a_j=\alpha_j(\mathbf{o}_j)$. Let $\rho$ be the action pattern specifying the
occurrence of $p$ in $a_i$, and $\sigma$ the pattern specifying its occurrence
in $a_j$. We write
$\mathbf{o}_i^\rho=(o_{i,\rho_1},\ldots,o_{i,\rho_r})$
for the object arguments selected by $\rho$. The head computes
\begin{align}\label{eq:scores}
S_{p,\rho,\sigma}(i,j)
&=
q_{p,\alpha_i,\rho}\,
k_{p,\alpha_j,\sigma}\,
\mathbb{I}\!\left[
\mathbf{o}_i^\rho=\mathbf{o}_j^\sigma
\right].
\end{align}
The $q$ term indicates whether the occurrence in $a_i$ is a precondition, the $k$ term whether the occurrence in $a_j$ is an effect, and the indicator ensures that both refer to the same ground instance of $p$. The equality test is implemented using one-hot object representations and structured query and key projections, as detailed in Appendix~\ref{app:strips-transformer}.

We then apply stick-breaking attention to these scores. For binary scores, it selects the most recent effect on the same ground instance of $p$, while $v$ indicates whether this effect deletes $p$. The resulting output $Y_{p,\rho}(i)$ is therefore $1$ exactly when the corresponding precondition is false.

The precondition outputs are combined as a continuous OR,
\begin{align}
Y(i)
&=
1-\prod_{p\in\mathcal{P}}\prod_{\rho}
\left(1-Y_{p,\rho}(i)\right).
\end{align}
For binary outputs, $Y(i)=1$ exactly when at least one precondition is false,
and hence when $a_i$ is inapplicable.
The full attention computation is given in
Appendix~\ref{app:strips-transformer}.

The same computation can incorporate state information by encoding it as auxiliary actions with fixed parameters. Before the domain trace, the action $\texttt{init-start}$ sets all ground predicate instances to false, and $\texttt{init-}p(\mathbf{o})$ actions set those in the initial state to true. For final-state supervision, $\texttt{test-}p(\mathbf{o})$ is appended after the trace, with $p(\mathbf{o})$ as its only precondition and no effects. Its predicted applicability therefore indicates whether $p(\mathbf{o})$ holds in the resulting state.

Finally, the learned parameters are indexed by predicates, action schemas,
and action patterns, but not by object identities or by the number $N$ of
objects. Changing $N$ only changes the fixed object representations used to
implement the equality test above. The same parameters can therefore be
instantiated on problems with a different number of objects. After
binarization, they directly specify the learned lifted \strips model.

Appendix~\ref{app:brap-program} describes the corresponding B-RASP program for the LS Transformer, which extends propositional formulation of~\citet{nunez2026strips}.